\begin{afterword}
\summaryhead

The post opens with a passage from Dale Griffin and Amos Tversky. They asked colleagues choosing between job offers how likely they were to choose each job. The average confidence in the predicted choice was 66\%, yet only 1 of the 24 chose the other job.

The author reports that reading this changed the way they thought. Once you can guess which way you will decide, you have in all probability already decided, and most of the time you can guess within half a second of hearing the question. The moment before you can guess is short, and it is the only time in which intelligence can act, for beliefs as for choices.

So the plan of reaching a belief by any means and then rejecting it if it cannot be justified will not work, because we change our minds much less often than we think. We remember the times we changed our minds more easily than the times we did not, and many biases keep an idea in place once it is in your head.

\responsehead

The post's evidence is one paragraph from Griffin and Tversky, and the post quotes it faithfully. The survey supports the title in a narrow sense: people switched from their predicted choice less often than their stated confidence implied. Had the 24 colleagues been well calibrated, about eight of them would have chosen the job they thought less likely. One did.

The post draws more from the survey than it shows. Griffin and Tversky present the result as underconfidence. They explain it by confidence that follows the balance of arguments: a person with a 2-to-1 balance of reasons feels about two-thirds sure, although even a small advantage for one job ``would generally lead to the choice'' of it. On that account no one needs to have decided early. The survey measured how well people predicted their choice. It did not measure when they decided, or whether a new reason could have changed the choice. The post's reading, that once you can guess your answer you have ``in all probability, already decided,'' is its own inference, and the half second has no source.

The general form of the title also goes beyond the source. On the same page Griffin and Tversky note that people are sometimes overconfident in predicting their own behavior, and cite a study in which students' predictions of their own actions were consistently overconfident. In the paper, which way the error runs depends on the evidence a person has. One survey of job choices cannot settle it in general, and the post extends it to ``questions of fact'' as well.

The step to the post's target is the least supported. The post argues against reaching a belief by non-rational means and then checking whether it can be justified. The survey concerns people weighing two job offers. It does not observe anyone checking a conclusion and failing to drop it. Evidence closer to that case exists in the studies of belief perseverance, where impressions outlast the discrediting of their source, and it supports the post's closing sentence. The post offers a list of bias names instead.

In short: a faithfully quoted survey of 24 job choices, read as showing that people decide early and rarely change their minds, although its authors reported it as underconfidence and noted that people sometimes overrate how well they can predict themselves.
\end{afterword}
